\chapter{Graph Neural Networks}
\label{chap:graph_neural_networks}

\begin{tldr}
Graph neural networks extend deep learning to non-Euclidean data---social networks, molecules, transaction graphs, knowledge bases---where entities are nodes and relationships are edges. The core idea is simple: iteratively update each node's representation by aggregating information from its neighbors (message passing). This chapter covers when to reach for GNNs instead of tabular or sequence models, the key architectures (GCN, GraphSAGE, GAT, GIN) and their trade-offs, the critical challenge of over-smoothing, scalability techniques for billion-node graphs, and practical applications. GNNs appear less frequently in interviews than transformers but are a strong differentiator for Staff+ candidates---especially for fraud detection, recommendation, and molecular design questions.
\end{tldr}

%==============================================================================
\section{When to Use GNNs}
\label{sec:when_gnns}
%==============================================================================

Not every relational dataset needs a GNN. Most candidates reach for GNNs too eagerly. The first question should always be: \textit{does the graph structure carry signal that a flat model would miss?}

\subsection{Graph Representation Basics}

A graph $G = (V, E)$ consists of nodes $V$ and edges $E$. Each node $v$ may have features $\vx_v \in \R^d$. Edges can be directed or undirected, weighted or unweighted, and may carry their own features. The structure is typically represented as an adjacency matrix $\mA \in \R^{|V| \times |V|}$ or, more practically, as edge lists.

\textbf{Common graph types in ML:}
\begin{itemize}
    \item \textbf{Homogeneous:} One node type, one edge type (social follow graph)
    \item \textbf{Heterogeneous:} Multiple node types, multiple edge types (user-item-store graph in e-commerce)
    \item \textbf{Bipartite:} Two node types, edges only between types (user-item interactions)
    \item \textbf{Temporal:} Edges have timestamps (transaction networks, communication graphs)
    \item \textbf{Knowledge graph:} Directed, labeled edges as (head, relation, tail) triples
\end{itemize}

\subsection{Decision Criteria: GNN vs.\ Alternatives}

\begin{table}[H]
\centering
\caption{When to use GNNs vs.\ alternative approaches}
\begin{tabularx}{\textwidth}{lXX}
\toprule
\textbf{Signal Source} & \textbf{Approach} & \textbf{Example} \\
\midrule
Node features only & Standard ML (XGBoost, MLP) & Predict user churn from profile features \\
Graph topology only & Graph algorithms (PageRank, label propagation) & Identify hub nodes in social networks \\
Node features + local structure & GNN & Fraud detection (account features + transaction patterns) \\
Node features + global structure & GNN with positional encoding & Molecular property prediction \\
Relational data, no explicit graph & Build graph + GNN \textit{or} tabular & Depends on relational signal strength \\
Very sparse / disconnected graph & Tabular often wins & Most nodes have $<$2 edges \\
\bottomrule
\end{tabularx}
\end{table}

\begin{keyinsight}{GNNs Add Value When Neighbors Are Informative}
The fundamental assumption behind GNNs is \textit{homophily}: connected nodes tend to share labels or properties. A fraudster's neighbors are more likely to be fraudsters. A molecule's biological activity depends on how its atoms are bonded, not just which atoms are present. If this assumption holds, GNNs outperform flat models. If neighbors carry no signal (e.g., random edges), a GNN reduces to a noisy MLP. Always validate that graph structure improves over a node-features-only baseline before committing to a GNN.
\end{keyinsight}

\subsection{Common Task Types}

\begin{table}[H]
\centering
\caption{Graph ML task taxonomy}
\begin{tabularx}{\textwidth}{lXX}
\toprule
\textbf{Task} & \textbf{Definition} & \textbf{Applications} \\
\midrule
Node classification & Predict a label for each node & Fraud detection, user categorization \\
Link prediction & Predict whether an edge exists & Friend recommendation, drug-target interaction \\
Graph classification & Predict a label for the entire graph & Molecular property prediction, malware detection \\
Edge classification & Predict a label for each edge & Relation typing in knowledge graphs \\
Node clustering & Unsupervised grouping of nodes & Community detection, customer segmentation \\
\bottomrule
\end{tabularx}
\end{table}

%==============================================================================
\section{Key Architectures}
\label{sec:gnn_architectures}
%==============================================================================

\subsection{The Message Passing Framework}

Every major GNN architecture follows the same template. Understanding this template is more important than memorizing individual variants.

\begin{mathresult}{Message Passing Neural Network (MPNN)}
For each node $v$, at layer $l$:
\begin{equation}
\vh_v^{(l)} = \text{UPDATE}\!\left(\vh_v^{(l-1)},\; \text{AGG}\!\left(\left\{\,\text{MSG}\!\left(\vh_v^{(l-1)}, \vh_u^{(l-1)}, e_{vu}\right) : u \in \mathcal{N}(v)\,\right\}\right)\right)
\label{eq:mpnn}
\end{equation}
where $\mathcal{N}(v)$ is the neighbor set of $v$, and $\vh_v^{(0)} = \vx_v$ (input features).
\end{mathresult}

Three design choices define each architecture:
\begin{enumerate}
    \item \textbf{MSG} --- how to compute the message from each neighbor
    \item \textbf{AGG} --- how to aggregate messages (sum, mean, max, attention-weighted)
    \item \textbf{UPDATE} --- how to combine the aggregated message with the node's own representation
\end{enumerate}

The number of message passing layers $L$ determines how far information propagates: after $L$ layers, each node's representation incorporates information from its $L$-hop neighborhood. This is the GNN analogue of receptive field in CNNs.

\subsection{GCN: Graph Convolutional Network}

\textbf{What it is.} The foundational GNN architecture (Kipf and Welling, 2017). Uses a symmetric normalized aggregation that averages neighbor features, weighted by degree.

\begin{mathresult}{GCN Layer}
\begin{equation}
\vh_v^{(l)} = \sigma\!\left(\mW^{(l)} \cdot \text{MEAN}\!\left(\left\{\,\frac{\vh_u^{(l-1)}}{\sqrt{|\mathcal{N}(v)| \cdot |\mathcal{N}(u)|}} : u \in \mathcal{N}(v) \cup \{v\}\,\right\}\right)\right)
\label{eq:gcn}
\end{equation}
where $\sigma$ is a nonlinearity (ReLU), and the self-loop ($v \in \mathcal{N}(v)$) ensures the node retains its own information.
\end{mathresult}

\textbf{When to use.} Good default for homogeneous graphs with moderate size. Simple, fast, well-understood.

\textbf{Pitfalls.} Requires the full adjacency matrix at training time (transductive by default---cannot handle unseen nodes without retraining). The symmetric normalization treats all neighbors equally, which limits expressivity.

\subsection{GraphSAGE}

\textbf{What it is.} SAmple and aggreGatE (Hamilton et al., 2017). Designed for \textit{inductive} learning on large graphs. Two key differences from GCN: (1) samples a fixed-size neighborhood instead of using all neighbors, enabling mini-batch training; (2) separates the node's own features from the aggregated neighbor features with a concatenation.

\textbf{How it works:}
\begin{enumerate}
    \item Sample $k$ neighbors per node at each layer (e.g., $k_1 = 25$, $k_2 = 10$ for two layers)
    \item Aggregate sampled neighbor features (mean, LSTM, or max pooling)
    \item Concatenate the node's own embedding with the aggregation, then apply a linear transform
\end{enumerate}

\textbf{When to use.} Large graphs (millions to billions of nodes) where full-batch GCN is infeasible. Production systems where new nodes appear over time (inductive setting).

\textbf{Pitfalls.} Neighbor sampling introduces variance. The sampling budget creates a trade-off: larger samples are more accurate but increase the computation tree exponentially (with $L$ layers and $k$ samples per layer, the computation tree has $O(k^L)$ nodes).

\subsection{GAT: Graph Attention Network}

\textbf{What it is.} Applies attention over neighbors to learn \textit{which} neighbors matter most (Velickovic et al., 2018). Instead of treating all neighbors equally (GCN) or sampling uniformly (GraphSAGE), GAT computes a learned attention weight for each edge.

\begin{mathresult}{GAT Attention}
\begin{equation}
\alpha_{vu} = \frac{\exp\!\left(\text{LeakyReLU}\!\left(\va^T [\mW\vh_v \,\|\, \mW\vh_u]\right)\right)}{\sum_{w \in \mathcal{N}(v)} \exp\!\left(\text{LeakyReLU}\!\left(\va^T [\mW\vh_v \,\|\, \mW\vh_w]\right)\right)}
\label{eq:gat}
\end{equation}
where $\|$ is concatenation, $\va$ is a learnable attention vector, and $\mW$ projects node features.
\end{mathresult}

GAT uses multi-head attention (same concept as in Transformers): $H$ independent attention heads, with outputs concatenated (intermediate layers) or averaged (final layer).

\textbf{When to use.} Heterogeneous neighborhoods where neighbor importance varies significantly. Works well when different neighbors contribute different types of information.

\textbf{Pitfalls.} The attention mechanism in GATv1 is a function of individual node features, not their \textit{interaction}---it can rank neighbors the same way regardless of the target node. GATv2 (Brody et al., 2022) fixes this by applying the nonlinearity after concatenation, enabling truly dynamic attention. Attention also adds compute and memory cost per edge.

\subsection{Architecture Comparison}

\begin{table}[H]
\centering
\caption{GNN architecture comparison for interview reference}
\begin{tabularx}{\textwidth}{lXXXXc}
\toprule
\textbf{Model} & \textbf{Aggregation} & \textbf{Strengths} & \textbf{Weaknesses} & \textbf{Scale} & \textbf{Year} \\
\midrule
GCN & Normalized mean & Simple, fast, solid baseline & Transductive, equal neighbor weight & Medium & 2017 \\
GraphSAGE & Sampled + concat & Inductive, scalable & Sampling variance, exponential fan-out & Large & 2017 \\
GAT & Attention-weighted & Learns neighbor importance & Higher compute, GATv1 attention is limited & Medium & 2018 \\
GIN & Sum (injective) & Most expressive (WL-equivalent) & Can overfit on small graphs & Small--Med & 2019 \\
Graph Transformer & Full attention (all nodes) & Global context, no locality bias & Quadratic in nodes, needs positional encoding & Small--Med & 2020+ \\
\bottomrule
\end{tabularx}
\end{table}

\subsection{GIN and the Weisfeiler-Leman Connection}

\textbf{What it is.} The Graph Isomorphism Network (Xu et al., 2019) is designed to be \textit{maximally expressive} among message-passing GNNs. Its aggregation uses sum (not mean or max), which preserves multiset structure and makes the layer update injective.

\textbf{Why it matters (the WL test intuition).} The Weisfeiler-Leman (WL) graph isomorphism test is a classical algorithm that iteratively refines node labels by hashing each node's label together with the \textit{multiset} of its neighbors' labels. Xu et al.\ proved that message-passing GNNs are \textit{at most} as powerful as the 1-WL test at distinguishing non-isomorphic graphs. GIN achieves this upper bound; GCN and GraphSAGE (which use mean aggregation) do not---they cannot distinguish certain graph structures that sum aggregation can.

\textbf{Concrete example.} Mean aggregation cannot distinguish a node with two neighbors of feature $[1, 0]$ from a node with one such neighbor---both produce mean $[1, 0]$. Sum aggregation produces $[2, 0]$ vs.\ $[1, 0]$, preserving the count.

\textbf{When to use.} Graph classification tasks where structural discrimination matters (molecular property prediction, graph-level anomaly detection).

\textbf{Pitfalls.} Being theoretically most expressive does not always mean best performance. On node classification with noisy features, the additional expressivity can lead to overfitting. GIN's advantage is most visible on small graph classification benchmarks.

\begin{warningbox}
All standard message-passing GNNs are bounded by the 1-WL test. They cannot distinguish certain non-isomorphic graphs (e.g., specific regular graphs with the same degree sequence). If your task requires distinguishing such structures, you need higher-order GNNs, subgraph methods, or positional/structural encodings. This is a common interview follow-up: ``What are the theoretical limitations of message passing?''
\end{warningbox}

%==============================================================================
\section{Challenges}
\label{sec:gnn_challenges}
%==============================================================================

\subsection{Over-Smoothing}

\textbf{What it is.} As you stack more GNN layers, node representations converge to indistinguishable values. After $L$ layers of message passing, each node's representation is a weighted average of its $L$-hop neighborhood. For dense graphs, a few layers suffice to cover the entire graph, causing all node representations to collapse to a near-identical vector. This is the GNN equivalent of vanishing gradients---but worse, because it is a fundamental property of the aggregation, not just a training issue.

\textbf{Why it happens.} Message passing is a form of Laplacian smoothing. Each layer averages a node's features with its neighbors, which is equivalent to applying a low-pass filter on the graph. Repeated application of a low-pass filter kills all high-frequency (local, discriminative) signal, leaving only the graph's leading eigenvector.

\begin{keyinsight}{Over-Smoothing Limits GNN Depth in Practice}
Most GNNs work best with 2--3 layers. Beyond that, performance typically degrades. This is not because deeper GNNs cannot be trained (they can) but because the representations lose discriminative power. A 2-layer GNN covers the 2-hop neighborhood, which is sufficient for most tasks. Deeper GNNs are needed only when the task genuinely requires long-range graph dependencies---and in that case, you need explicit anti-smoothing mechanisms.
\end{keyinsight}

\begin{table}[H]
\centering
\caption{Over-smoothing solutions}
\begin{tabularx}{\textwidth}{lXX}
\toprule
\textbf{Solution} & \textbf{How It Works} & \textbf{Trade-off} \\
\midrule
Residual connections & Skip connections between GNN layers: $\vh^{(l)} = \vh^{(l-1)} + \text{GNN}(\vh^{(l-1)})$ & Simple; partially mitigates but does not fully solve \\
JKNet (Jumping Knowledge) & Concatenate or max-pool representations from all layers as final embedding & Increases output dimension; adds layer-selection mechanism \\
DropEdge & Randomly remove edges during training (analogous to dropout for graphs) & Reduces over-smoothing by limiting propagation; adds noise \\
PairNorm & Normalize node embeddings to maintain constant total pairwise distance & Prevents collapse; may reduce expressivity \\
Graph Transformer & Replace local message passing with global attention over all nodes & Avoids locality constraint; quadratic cost in $|V|$ \\
Subgraph sampling & Extract local subgraphs and process independently & Limits propagation depth; loses global context \\
\bottomrule
\end{tabularx}
\end{table}

\textbf{The practical rule.} Start with 2 layers. If the task needs longer-range information, add layers one at a time with residual connections and monitor node embedding similarity (e.g., mean cosine similarity between representations). If similarity approaches 1.0, you have over-smoothed.

\subsection{Over-Squashing}

\textbf{What it is.} The complementary problem to over-smoothing. When a node needs information from $L$ hops away, that information must be compressed through a fixed-width bottleneck at each intermediate node. Exponentially many distant nodes' messages get ``squashed'' into a fixed-dimensional vector. This is analogous to the information bottleneck in deep RNNs.

\textbf{Mitigation.} Graph rewiring (add ``shortcut'' edges to reduce diameter), virtual nodes (a global node connected to all others that acts as a broadcast channel), or Graph Transformers (bypass locality entirely with global attention).

\subsection{Scalability}

Real-world graphs have billions of nodes and edges. Full-batch training is impossible.

\begin{table}[H]
\centering
\caption{GNN scalability techniques}
\begin{tabularx}{\textwidth}{lXXX}
\toprule
\textbf{Method} & \textbf{How It Works} & \textbf{Pro} & \textbf{Con} \\
\midrule
Neighbor sampling (GraphSAGE) & Sample $k$ neighbors per node per layer & Inductive, simple & Exponential fan-out with depth \\
Cluster-GCN & Partition graph into clusters, train on subgraphs & Reduces fan-out & Inter-cluster edges lost per batch \\
GraphSAINT & Sample subgraphs with importance sampling & Unbiased gradient estimates & Complex sampling scheme \\
Mini-batch with historical embeddings & Cache neighbor embeddings from prior iterations & Reduces recomputation & Stale embeddings add noise \\
\bottomrule
\end{tabularx}
\end{table}

\textbf{Neighbor sampling fan-out math.} For a 2-layer GNN with fan-out $[25, 10]$, computing one target node's embedding requires sampling $25 \times 10 = 250$ second-hop nodes plus 25 first-hop nodes $= 275$ total. For a batch of 1024 target nodes, you process up to $1024 \times 275 \approx 280$K node features per batch. This is manageable but grows quickly with more layers or larger fan-out.

\textbf{Cluster-GCN} (Chiang et al., 2019) avoids the fan-out problem entirely by partitioning the graph (e.g., using METIS) and training on entire subgraphs. Within a subgraph, full-batch message passing is used. The cost: edges between clusters are dropped during training, introducing bias. Mixing clusters across mini-batches partially mitigates this.

\begin{warningbox}
The exponential fan-out problem in neighbor sampling means that a 3-layer GNN with fan-out $[25, 10, 10]$ touches $25 \times 10 \times 10 = 2{,}500$ nodes per target node. This limits practical GNN depth even when over-smoothing is not a concern. If your interviewer asks ``why not just add more layers to capture longer-range dependencies?'' the answer is both over-smoothing \textit{and} the computational cost of deep neighbor sampling.
\end{warningbox}

%==============================================================================
\section{Applications}
\label{sec:gnn_applications}
%==============================================================================

\subsection{Social Networks}

\textbf{Tasks:} Community detection, influence prediction, friend recommendation, content recommendation.

\textbf{Why GNNs help.} Social networks exhibit strong homophily (friends share interests) and structural roles (influencers have distinct neighborhoods). GNNs capture both through message passing. For friend recommendation, link prediction with GNN-generated embeddings outperforms collaborative filtering on cold-start users because the graph structure provides signal even without interaction history.

\textbf{Scale challenge.} Social graphs are among the largest (billions of nodes). Production systems at Meta use GraphSAGE-style sampling with feature caching and distributed training.

\subsection{Molecular Property Prediction}

\textbf{Tasks:} Predict solubility, toxicity, binding affinity, or other physicochemical properties from molecular structure.

\textbf{Why GNNs are natural.} A molecule is literally a graph: atoms are nodes, bonds are edges. Both node features (atom type, charge, hybridization) and edge features (bond type, stereochemistry) carry signal. Graph-level prediction (pooling all node embeddings into a single graph embedding) is the standard approach.

\textbf{Architecture choice.} GIN or SchNet (which incorporates 3D coordinates as continuous edge features) are common. Attention-based variants (e.g., GPS---General Powerful Scalable graph transformer) that combine message passing with global attention are state-of-the-art on benchmarks like OGB-LSC.

\textbf{Practical note.} Molecular graphs are small (typically 10--100 nodes), so scalability is not a concern. The challenge is data efficiency---labeled molecular data is expensive to obtain (wet lab experiments), making few-shot learning and pretraining (e.g., self-supervised on unlabeled molecules) important.

\subsection{Fraud Detection}

\textbf{Tasks:} Identify fraudulent accounts, transactions, or reviews in a transaction or interaction graph.

\textbf{Why GNNs help.} Fraudsters create networks of colluding accounts. Individual features (e.g., account age, transaction frequency) can be spoofed, but the graph structure---dense clusters of new accounts transacting with each other in unusual patterns---is much harder to fake. GNNs detect these structural anomalies by propagating neighborhood information.

\textbf{Key design decisions:}
\begin{itemize}
    \item \textbf{Graph construction:} How to define edges (co-transaction, shared device, shared IP, shared payment method). Multiple edge types create a heterogeneous graph.
    \item \textbf{Label scarcity:} Fraud labels are rare ($<$1\% positive rate). Combine GNN with oversampling, focal loss, or semi-supervised approaches.
    \item \textbf{Temporal dynamics:} Fraud patterns evolve. Use temporal edges and retrain frequently, or use dynamic GNNs that model edge timestamps.
    \item \textbf{Adversarial robustness:} Fraudsters adapt. The model must be robust to adversarial graph manipulation (e.g., creating benign-looking edges to evade detection).
\end{itemize}

\subsection{Knowledge Graphs}

\textbf{Tasks:} Link prediction (knowledge base completion), entity classification, relation extraction.

\textbf{Representation.} Knowledge graphs store facts as (head, relation, tail) triples, e.g., (Einstein, born\_in, Germany). The graph is directed and multi-relational.

\textbf{Approaches:}
\begin{itemize}
    \item \textbf{Embedding methods} (TransE, RotatE, ComplEx): Learn entity and relation embeddings such that $\vh + \mathbf{r} \approx \mathbf{t}$ for valid triples. Simple, scalable, but do not use node features.
    \item \textbf{GNN-based} (R-GCN, CompGCN): Apply message passing with relation-specific transformations. Each relation type has its own weight matrix (or a basis decomposition to keep parameters manageable). These methods can incorporate node features and multi-hop reasoning.
\end{itemize}

\textbf{Practical note.} Knowledge graphs in industry (Google Knowledge Graph, Wikidata) have billions of triples. Embedding methods dominate at scale due to simplicity. GNN-based methods are preferred when node features are available or multi-hop reasoning is needed.

\subsection{Application Selection Guide}

\begin{table}[H]
\centering
\caption{GNN application quick reference}
\begin{tabularx}{\textwidth}{lXXl}
\toprule
\textbf{Domain} & \textbf{Task} & \textbf{Recommended Architecture} & \textbf{Graph Size} \\
\midrule
Social networks & Node class.\ / link pred. & GraphSAGE (scalable, inductive) & Billions \\
Molecules & Graph classification & GIN, SchNet, GPS & Small \\
Fraud detection & Node classification & R-GCN (heterogeneous), GraphSAGE & Millions \\
Knowledge graphs & Link prediction & R-GCN, TransE/RotatE & Billions \\
Recommendation & Link prediction & LightGCN, PinSage & Billions \\
Combinatorial optim. & Graph-level regression & GNN + set pooling & Small--Med \\
\bottomrule
\end{tabularx}
\end{table}

%==============================================================================
\section{Interview Questions}
\label{sec:gnn_interview}
%==============================================================================

\begin{interviewq}
\textbf{Q: Design a fraud detection system using GNNs.}

\textbf{A:}

\textbf{Step 1: Problem formulation.} Binary node classification: label each account as fraudulent or legitimate. The key insight is that fraud is a \textit{network} phenomenon---fraudulent accounts cluster together and exhibit unusual connectivity patterns that are invisible to per-account feature models.

\textbf{Step 2: Graph construction.} Build a heterogeneous multi-relational graph:
\begin{itemize}
    \item \textbf{Nodes:} User accounts, devices, IP addresses, payment instruments
    \item \textbf{Edges:} User-device (login), user-IP (session), user-payment (owns), user-user (transaction), user-user (shared device within 24h)
    \item Edge types are critical---shared-device edges between unrelated accounts are the strongest fraud signal.
\end{itemize}

\textbf{Step 3: Model.} Two-layer R-GCN (relational GCN) with relation-specific weight matrices (using basis decomposition to reduce parameters). Use GraphSAGE-style neighbor sampling for scalability (sample 15 neighbors per edge type per layer). Output: per-node fraud probability via sigmoid on the final embedding.

\textbf{Step 4: Handling class imbalance.} Fraud rate is typically 0.1--1\%. Use focal loss ($\gamma = 2$) or a weighted binary cross-entropy with upweight on positives. At evaluation, optimize for precision at fixed recall (e.g., 90\% recall) rather than accuracy or AUC alone, since false positives have a direct user-experience cost.

\textbf{Step 5: Serving and temporal aspects.} Fraud detection is real-time. Precompute node embeddings periodically (e.g., hourly) and cache them. For a new transaction, look up the sender and receiver embeddings, concatenate with transaction features, and run through a lightweight classifier. Retrain the GNN daily on the most recent labeled data, including investigator-confirmed fraud labels with a 24--48h delay.

\textbf{Step 6: Monitoring.} Track precision@90\%-recall over time. Monitor graph statistics (average degree, new edge rate) for distribution shift. A sudden spike in new accounts creating dense subgraphs is itself a fraud signal.

\textbf{What the interviewer is testing:} Can you design a complete ML system (not just pick a model)? Do you understand why GNNs specifically add value here (structural signal that per-account features miss)? Do you handle the practical challenges: class imbalance, real-time serving, temporal dynamics, and adversarial adaptation?

\textbf{Common follow-ups:}
\begin{itemize}
    \item ``How do you handle adversarial adaptation---fraudsters learning to avoid detection?'' (Feature rotation, ensemble with rule-based signals, red-team the model regularly, avoid publishing detection thresholds.)
    \item ``What if GNN training is too slow for hourly updates?'' (Use incremental graph updates with cached embeddings; only recompute for changed neighborhoods.)
    \item ``How do you handle cold-start accounts with no edges?'' (Fall back to the node-features-only MLP; the GNN provides no additional signal without edges.)
    \item ``How would you A/B test this?'' (Carefully---you cannot split users independently because fraud is a network effect. Use time-based switchbacks or measure on held-out investigator queues.)
\end{itemize}
\end{interviewq}

\begin{interviewq}
\textbf{Q: When would you use a GNN over tabular models on relational data?}

\textbf{A:}

The decision depends on whether the relational structure carries predictive signal \textit{beyond} what engineered features can capture.

\textbf{Use tabular models (XGBoost, LightGBM) when:}
\begin{itemize}
    \item Relational patterns can be captured by hand-crafted graph features: degree, PageRank, clustering coefficient, triangle count. These features fed into a gradient-boosted tree are often competitive with a GNN and much simpler to deploy.
    \item The graph is sparse or disconnected. If most nodes have fewer than 2 edges, there is minimal neighborhood information for a GNN to leverage.
    \item Deployment constraints favor simplicity. Tabular models are faster to train, easier to serve, simpler to debug, and do not require graph storage infrastructure.
    \item You have strong node-level features. If individual features are already highly predictive, the marginal gain from graph structure may not justify the complexity.
\end{itemize}

\textbf{Use GNNs when:}
\begin{itemize}
    \item The graph is dense and exhibits homophily---neighbors are informative about the target node. Validate this by checking if a simple label propagation baseline beats the features-only model.
    \item Hand-engineered graph features cannot capture the signal. Multi-hop patterns, subgraph motifs, and heterogeneous edge types are difficult to manually feature-engineer but are naturally captured by GNNs.
    \item The graph evolves and new nodes appear frequently (inductive GNNs like GraphSAGE handle this; static graph features must be recomputed).
    \item You need embeddings for downstream tasks (similarity search, clustering) where GNN embeddings capture richer structural information than manually crafted features.
\end{itemize}

\textbf{The pragmatic approach in practice:}
\begin{enumerate}
    \item Start with a tabular model on node features + hand-crafted graph statistics (baseline).
    \item Add a simple GNN (2-layer GCN) and compare.
    \item If the GNN improves by $>$2\% on the target metric, the structural signal is real. Invest in a more sophisticated GNN and proper graph infrastructure.
    \item If not, keep the tabular model. The engineering cost of GNNs (graph construction, neighbor sampling, serving) is only justified when the accuracy gain is substantial.
\end{enumerate}

\textbf{What the interviewer is testing:} Engineering maturity---do you default to the simplest approach and add complexity only when justified? Do you recognize that GNNs have significant infrastructure overhead? Can you articulate the specific conditions under which graph structure adds signal?

\textbf{Common follow-ups:}
\begin{itemize}
    \item ``How would you create graph features for the tabular baseline?'' (Degree, PageRank, ego-net density, neighbor label distribution, graph-based embeddings like node2vec.)
    \item ``What if the tabular model uses XGBoost and already has graph features---can a GNN still help?'' (Yes, if the signal is in multi-hop or structural patterns that are hard to hand-engineer. GNNs learn task-specific aggregation; manual features are generic.)
    \item ``How do you handle heterogeneous edges in the tabular approach vs.\ GNN?'' (Tabular: separate features per edge type. GNN: relation-specific transformations like R-GCN.)
\end{itemize}
\end{interviewq}

\begin{interviewq}
\textbf{Q: Explain over-smoothing in GNNs and how to mitigate it.}

\textbf{A:}

Over-smoothing is the phenomenon where stacking more GNN layers causes all node representations to converge to the same vector, making nodes indistinguishable.

\textbf{Why it happens.} Each GNN layer aggregates information from a node's neighbors. This is mathematically equivalent to applying a low-pass filter on the graph signal. After $L$ layers, each node's representation is a weighted combination of features from its entire $L$-hop neighborhood. For connected graphs with small diameter (most real-world graphs have diameter $<$10), even 3--4 layers mean each node has aggregated information from nearly the entire graph. The result: all nodes converge toward the graph's mean feature vector, scaled by the leading eigenvector of the normalized adjacency matrix.

\textbf{Quantifying it.} You can measure over-smoothing by tracking the mean pairwise cosine similarity of node embeddings across layers. In a healthy GNN (2 layers), this stays below 0.5. At 8+ layers on a dense graph, it approaches 1.0.

\textbf{Mitigation strategies:}
\begin{enumerate}
    \item \textbf{Keep it shallow.} The simplest and most effective approach. Most tasks do not need $>$3 hops of information. Use 2 layers as the default.
    \item \textbf{Residual connections.} Add skip connections: $\vh^{(l)} = \vh^{(l-1)} + \text{GNN}(\vh^{(l-1)})$. Preserves the original signal at each layer.
    \item \textbf{Jumping Knowledge Networks.} Concatenate representations from all layers $[\vh^{(1)} \| \vh^{(2)} \| \cdots \| \vh^{(L)}]$ and learn which layer's representation matters most per node via attention or max-pooling. Lets different nodes use different effective depths.
    \item \textbf{DropEdge.} Randomly remove a fraction of edges during training. This slows down information propagation and acts as regularization against over-smoothing.
    \item \textbf{Graph Transformers.} Replace local message passing with global attention. Since every node can attend to every other node directly, you do not need to stack layers for long-range information. But this is $O(|V|^2)$, so it only works for smaller graphs.
\end{enumerate}

\textbf{The deeper issue.} Over-smoothing is a symptom of a more fundamental problem: message-passing GNNs have limited expressive power. They are bounded by the 1-WL test, and depth does not overcome this limitation---it only makes it worse by destroying local signal. For tasks requiring global graph-level reasoning, the solution is not deeper GNNs but architectural changes: positional encodings, virtual nodes, or Graph Transformers.

\textbf{What the interviewer is testing:} Do you understand over-smoothing at a mechanistic level (Laplacian smoothing, low-pass filter), not just as a buzzword? Can you propose multiple solutions and explain their trade-offs? Do you know the practical implication (most GNNs should be 2--3 layers)?

\textbf{Common follow-ups:}
\begin{itemize}
    \item ``How is over-smoothing different from vanishing gradients?'' (Vanishing gradients is a training optimization issue solvable with better initialization and normalization. Over-smoothing is a \textit{representation} issue---even with perfect gradient flow, the representations converge. You can train a 10-layer GNN without gradient issues and still observe over-smoothing.)
    \item ``What is over-squashing and how does it relate?'' (Over-squashing is the complementary problem: messages from distant nodes get compressed through fixed-width bottlenecks at intermediate nodes, losing information. Over-smoothing kills discrimination; over-squashing kills long-range signal.)
    \item ``Would batch normalization help with over-smoothing?'' (Standard batch norm does not address it because it normalizes per feature, not per node pair. PairNorm, which normalizes to maintain inter-node distance, is specifically designed for this.)
\end{itemize}
\end{interviewq}
